\chapter{Claim: Valuations Predict the Equity Premium}
\label{ch:07}

This is the first of the book's claim chapters. Chapters~\ref{ch:01} to~\ref{ch:06} built the tools: statistics for noisy returns, corrections for searching, honest out-of-sample\index{in-sample versus out-of-sample} tests and trading costs\index{transaction costs}. From here on, each chapter takes one widely held claim and puts it on trial with those tools, following the same steps: how the idea works, the claim stated so it can fail, the evidence that made it famous, a replication\index{replication} on open data, stress tests, and a dated verdict.

The first claim is one of the oldest and most consequential in finance. If stock returns are predictable\index{return predictability} from valuations, then investors can time the market, retirement plans should assume lower returns when prices are high, and the efficient-markets view that dominated the 1970s needs rethinking. Few empirical questions have been argued over longer, by more distinguished economists, with more at stake for ordinary savers.

The chapter's answer depends on what ``predictable'' is supposed to mean. For deciding whether to be in or out of stocks next month or next year, the evidence is weak. For setting a realistic range for the next decade's returns, valuations carry real information, though the range is wide. The lab lets you see both conclusions emerge from the same data.

By the end of the chapter you will be able to:

\begin{itemize}
  \item explain the motivation behind return predictability research and its implications for financial planning, asset allocation, and economic theory;
  \item identify and apply the evaluation tools drawn from Chapters~\ref{ch:01} to~\ref{ch:06} specifically for testing return predictability;
  \item explain how return predictability connects to out-of-sample evaluation and differs from cross-sectional multi-factor models;
  \item compute present values\index{present value} and use the Gordon growth model\index{dividend discount model} to link a dividend yield\index{valuation ratios!dividend yield} to expected returns\index{return!expected} and growth;
  \item distinguish expected from realized returns\index{return!realized}, and real from nominal returns\index{real and nominal returns};
  \item read a yield curve\index{yield curve} and compute a term spread\index{term spread};
  \item explain what valuation ratios such as the dividend--price ratio and CAPE\index{valuation ratios!CAPE} measure, and why the Campbell--Shiller\index{Campbell--Shiller decomposition} identity says they must carry information about the future;
  \item state the claim that returns are predictable in a form that can fail, and summarize the evidence for and against it;
  \item test the claim with corrected long-horizon regressions, a bootstrap\index{bootstrap} for the Stambaugh\index{Stambaugh bias} bias and walk-forward\index{walk-forward evaluation} forecasts;
  \item use valuations to set a range for long-run returns, rather than to time the market.
\end{itemize}

\section{Motivation: why this claim matters and its implications}
\label{sec:07-motivation}

Return predictability sits at the intersection of economic theory, institutional asset allocation, and personal wealth management. Whether equity returns can be forecast from aggregate valuation ratios is not merely an academic dispute; it alters how trillions of dollars in capital are managed and evaluated.

\subsection{Why researchers and practitioners care}

\begin{itemize}
  \item \textbf{Tactical asset allocation and market timing:} If aggregate valuation ratios (such as the dividend yield or CAPE) contain reliable predictive signals, investors can systematically shift between equities and risk-free assets to improve risk-adjusted returns or reduce drawdown risk.
  \item \textbf{Long-term wealth and retirement planning:} Pension plans, university endowments, and individual retirees rely on return assumptions over multi-year or multi-decade horizons. Assuming a static historical mean return (e.g., 7\% real) regardless of starting valuation levels can lead to severe underfunding when markets are historically expensive.
  \item \textbf{Foundation of modern asset pricing theory:} Economists care because return predictability directly tests the Efficient Market Hypothesis (EMH). If market returns vary over time, asset pricing models must explain whether this predictability reflects rational, time-varying compensation for risk or persistent market mispricing driven by investor psychology.
\end{itemize}

\subsection{Key implications}

The empirical validity of return predictability carries distinct implications across finance:

\begin{itemize}
  \item \textbf{For investors and asset managers:} If short-horizon predictability fails out of sample, tactical market timing destroys value through trading drag and missed market gains. Conversely, if long-horizon predictability survives weakly, valuations should inform baseline return expectations rather than dynamic timing rules.
  \item \textbf{For corporate finance and cost of capital:} Capital expenditure decisions and equity issuance timing depend on the cost of equity. If market expected returns fluctuate with valuations, a firm's hurdle rate must adjust over the business cycle rather than remain fixed.
  \item \textbf{For macroeconomic policy and financial stability:} Central banks and regulators monitor valuation metrics (such as the term spread or CAPE) to assess financial stability, credit expansion, and asset price bubble risks.
\end{itemize}

\section{Toolkit needed from earlier chapters}
\label{sec:07-connection}

As the first claim chapter, Chapter~\ref{ch:07} relies heavily on the methodological foundation laid in Chapters~\ref{ch:01} to~\ref{ch:06}. Understanding how this chapter connects to out-of-sample evaluation, how it differs from multi-factor models, and how it deploys tools from earlier chapters clarifies its position in the book.

\subsection{Relationship to Out-of-Sample Evaluation (Chapter~\ref{ch:05})}

Chapter~\ref{ch:07} is the direct empirical application of the framework built in Chapter~\ref{ch:05}:

\begin{itemize}
  \item \textbf{Same core task (time-series market prediction):} Both chapters evaluate predictive regressions of the aggregate equity premium over time ($r_{t+1} = \alpha + \beta x_t + \varepsilon_{t+1}$).
  \item \textbf{Shared toolkit and data:} Chapter~\ref{ch:07} uses the exact same Goyal and Welch monthly dataset and out-of-sample testing design established in Chapter~\ref{ch:05} (walk-forward windows, purging overlapping targets, out-of-sample $R^2$, and the Clark--West test).
  \item \textbf{Key difference in focus:} Chapter~\ref{ch:05} establishes the evaluation methodology---building loss functions, establishing the historical mean benchmark, identifying leakage, and defining statistical tests for model comparison. Chapter~\ref{ch:07} puts a specific economic claim on trial---testing whether valuation ratios (dividend--price ratio, CAPE) and yield spreads carry genuine predictive power for aggregate market returns across different time horizons (1-month, 1-year, and 10-year), backed by present-value theory (the Campbell--Shiller identity).
\end{itemize}

\subsection{Difference from multi-factor models}

Chapter~\ref{ch:07} differs fundamentally from multi-factor models in both economic scope and econometric structure:

\begin{itemize}
  \item \textbf{Time-series versus cross-section:} Chapter~\ref{ch:07} is a time-series problem: it asks whether a single asset class (the broad stock market) can be timed over time using valuation signals. Multi-factor models are primarily a cross-sectional problem: they explain why different assets or portfolios (e.g., small-cap vs.\ large-cap, value vs.\ growth) earn different average returns over the same period based on their risk factor exposures ($\beta$'s).
  \item \textbf{Predictors versus risk factors:} Chapter~\ref{ch:07} uses aggregate macroeconomic and valuation indicators (CAPE, dividend yield, term spread) to forecast market returns. Multi-factor models use systematic risk factors (such as market excess return, SMB, HML, or RMW) to decompose portfolio risk and test for abnormal return ($\alpha$).
\end{itemize}

\subsection{Summary of tools drawn from Chapters 1--6}

Table~\ref{tab:07-tools} explicitly identifies the specific analytical tools drawn from Chapters~\ref{ch:01} through~\ref{ch:06} that are deployed throughout this chapter to evaluate market predictability.

\begin{table}[htbp]
\centering
\small
\caption{Tools drawn from Chapters 1--6 used in Chapter 7.}
\label{tab:07-tools}
\begin{tabularx}{0.95\linewidth}{@{}>{\raggedright\arraybackslash}p{0.18\linewidth}>{\raggedright\arraybackslash}p{0.28\linewidth}>{\raggedright\arraybackslash}X@{}}
\toprule
Origin Chapter & Tool / Concept & Role in Chapter 7 \\
\midrule
Chapter~\ref{ch:01} & Expected vs.\ Realized Returns & Separating predictable expected returns from large unforecastable return shocks. \\
Chapter~\ref{ch:02} & Point-in-Time Data \& Publication Lags & Lagging inflation and accounting earnings to eliminate look-ahead leakage in real-time evaluations. \\
Chapter~\ref{ch:03} & Overlapping Data Corrections (Newey--West, Hodrick) & Correcting standard errors when forecasting 1-year and 10-year returns sampled at monthly frequency. \\
Chapter~\ref{ch:04} & Search Hazards \& Multiple Testing & Accounting for data-mining bias when selecting among candidate valuation ratios and horizons. \\
Chapter~\ref{ch:05} & Walk-Forward Evaluation \& Purging & Generating genuine recursive forecasts without future leakage and purging overlapping targets. \\
Chapter~\ref{ch:05} & Out-of-Sample $R^2$ \& Clark--West Test & Measuring whether valuation forecasts beat the historical mean benchmark out of sample. \\
Chapter~\ref{ch:05} & Campbell--Thompson Restrictions & Applying sign and non-negativity economic bounds to prevent noisy slopes from destroying forecasts. \\
Chapter~\ref{ch:06} & Transaction Costs \& Turnover & Evaluating whether tactical market-timing rules generate profits after realistic execution drag. \\
\bottomrule
\end{tabularx}
\end{table}

\section{Foundations: valuing a stock market, interest rates and inflation}
\label{sec:07-foundations}

\subsection{Present value}

A dollar received in the future is worth less than a dollar today, because today's dollar could be invested. The rate used to translate future money into today's money is the \textbf{discount rate}\index{discount rate|textbf}. At a discount rate $r$ per year, a cash flow $C$ received in $n$ years has a \textbf{present value}\index{present value|textbf} of $C/(1 + r)^n$: \$110 received in one year is worth \$100 today at a 10\% discount rate. The price of any asset can be written as the present value of the cash flows it is expected to pay, discounted at the return investors require to hold it.

\subsection{The dividend discount model}

For a stock, or for the market as a whole, the cash flows are the dividends that shareholders receive. The model becomes simple if we make two assumptions. First, dividends grow forever at a constant rate $g$ each year. Second, investors require an expected return $r$ that is higher than $g$. Under these assumptions the present value of all future dividends reduces to a single fraction, known as the \textbf{Gordon growth model}\index{dividend discount model|textbf}:
\begin{equation}
  P_0 = \frac{D_1}{r - g}
  \label{eq:07-gordon}
\end{equation}
Here $P_0$ is today's price and $D_1$ is the dividend expected next year. In words, the price equals next year's dividend divided by the gap between the return investors require and the rate at which dividends grow. For example, a market that will pay a \$3 dividend next year, with dividends growing at 4\% a year and a required return of 7\%, is worth $\$3/(0.07 - 0.04) = \$100$.

Dividing both sides of Equation~\eqref{eq:07-gordon} by the price turns the model around, and this second form is the key to the chapter:
\begin{equation}
  \frac{D_1}{P_0} = r - g
  \label{eq:07-gordon-yield}
\end{equation}
The left-hand side is the \textbf{dividend yield}\index{valuation ratios!dividend yield|textbf}, so the equation says that the dividend yield equals the required return minus expected growth. In the example, the yield of $\$3/\$100 = 3\%$ is exactly the 7\% required return minus the 4\% growth rate. A high dividend yield must therefore mean that investors require high returns, that they expect low growth, or both. The same logic applies to earnings. If a firm pays out a constant share of its earnings as dividends, called the \textbf{payout ratio}\index{payout ratio|textbf}, then the \textbf{earnings yield}\index{valuation ratios!earnings yield|textbf} $E/P$, which is the inverse of the P/E ratio, is also tied to required returns and growth.

Because $r - g$ is small, the price is very sensitive to it (Figure~\ref{fig:07-gordon}). A one-point fall in the required return can raise the price by half. This is why stock markets can swing widely on changes in expected returns alone.

\providecolor{fenavy}{HTML}{1B2A6B}
\providecolor{feorange}{HTML}{B8481C}

\begin{figure}[htbp]
  \centering
  \begin{tikzpicture}[
    declare function={price(\x) = 3/((\x - 4)/100);}
  ]
  \begin{axis}[
    width=0.6\linewidth, height=0.38\linewidth,
    axis lines=left,
    axis line style={-stealth, thin},
    xmin=5, xmax=10.3, ymin=0, ymax=230,
    xtick={5,6,7,8,9,10},
    xticklabels={5\%,6\%,7\%,8\%,9\%,10\%},
    ytick={0,50,100,150,200},
    yticklabels={\$0,\$50,\$100,\$150,\$200},
    tick label style={font=\small},
    xlabel={Required return, $r$},
    ylabel={Price, $P = D/(r - g)$},
    samples=200, domain=5.35:10,
    clip=false,
  ]
    \addplot[fenavy, very thick] {price(x)};

    % r = 7%: price $100, yield 3%
    \draw[gray, densely dashed, thin] (axis cs:7,0) -- (axis cs:7,100) -- (axis cs:5,100);
    \fill[fenavy] (axis cs:7,100) circle (2.2pt);
    \node[font=\small, anchor=south west, align=left] at (axis cs:7.1,102)
      {$r = 7\%$: price \$100,\\dividend yield 3\%};

    % r = 6%: price $150, yield 2%
    \draw[feorange, densely dashed, thin] (axis cs:6,0) -- (axis cs:6,150) -- (axis cs:5,150);
    \fill[feorange] (axis cs:6,150) circle (2.2pt);
    \node[feorange, font=\small, anchor=south west, align=left] at (axis cs:6.1,152)
      {$r = 6\%$: price \$150,\\dividend yield 2\%};

    \node[font=\small, anchor=north east, align=right, text=gray!40!black]
      at (axis cs:10.2,215) {$D = \$3$, $g = 4\%$};
  \end{axis}
  \end{tikzpicture}
  \caption{The Gordon growth model for a market paying a \$3 dividend next
    year, with dividends growing 4\% a year. Lowering the required return from
    7\% to 6\% raises the price from \$100 to \$150 and lowers the dividend
    yield from 3\% to 2\%. High prices and low yields go with low expected
    returns.}
  \label{fig:07-gordon}
\end{figure}

\subsection{Expected and realized returns}

The required return $r$ in these formulas is the expected return of Section~\ref{sec:01-expret}: what investors anticipate, on average, when they buy. The \textbf{realized return} in any year is the expected return plus a surprise, and Section~\ref{sec:01-foundations} showed how large those surprises are for stocks. The same distinction applies to the equity premium (Section~\ref{sec:01-rf}): U.S. stocks have earned a realized premium of roughly 6 to 8 percent a year over Treasury bills, but how much of that was expected is debated.

\subsection{Interest rates and the term structure}
\label{sec:07-rates}

The \textbf{yield} on a bond is the single discount rate that makes the present value of its promised payments equal to its price. Yields on government bonds of different maturities form the \textbf{term structure}, or \textbf{yield curve} (Figure~\ref{fig:07-curve}). The \textbf{term spread} is the yield on a long-term bond (say, 10 years) minus the yield on a short-term bill. The curve normally slopes upward; it tends to be steep near the bottom of recessions and flat or \textbf{inverted}, with short rates above long ones, before them. Hull's \emph{Options, Futures, and Other Derivatives} covers the term structure in detail; Chapter~\ref{ch:16} tests what forward rates predict.

\subsection{Real and nominal}

\textbf{Nominal} returns and yields are in dollars; \textbf{real} returns subtract inflation, approximately $r_{\text{real}} \approx r_{\text{nominal}} - \text{inflation}$. A 5\% nominal return with 3\% inflation is about a 2\% real return. Stocks are claims on real businesses whose earnings tend to rise with prices over time, so their earnings yield behaves more like a real return. Government bond yields are nominal. Comparing the two directly, as the ``Fed model'' does, mixes the units.

\textbf{Worked example.} A market pays a dividend of \$3 next year on a price of \$100, a dividend yield of 3\%. If dividends are expected to grow 4\% a year, the Gordon model implies an expected return of $3\% + 4\% = 7\%$. If the price rose to \$150 with no change in expected growth, the yield would fall to 2\% and the implied expected return to 6\%.

\section{How it works}
\label{sec:07-how}

\textbf{Valuation ratios.} A valuation ratio compares a stock market's price with a measure of what it pays or earns.

\begin{itemize}
  \item The \textbf{dividend--price ratio} divides the last twelve months of dividends by the index level.
  \item The \textbf{cyclically adjusted price--earnings ratio (CAPE)}, introduced by Campbell and Shiller (1988, 1998), divides the inflation-adjusted price by the average of inflation-adjusted earnings over the previous ten years. The ten-year average smooths out the swings of the \textbf{business cycle}, the alternation of expansions and recessions.
\end{itemize}

When prices are high relative to dividends or earnings, the market is said to be expensive.

\textbf{Why a valuation ratio should carry information.} The present-value identity of Campbell and Shiller (1988) links today's dividend--price ratio to the future. Approximately, and without any assumption about investor behavior,
\begin{equation}
  d_t - p_t \approx \text{constant} + \sum_{j=1}^{\infty} \rho^{j-1}\left(r_{t+j} - \Delta d_{t+j}\right)
  \label{eq:07-identity}
\end{equation}
where $d - p$ is the log dividend--price ratio, $r$ is the log return, $\Delta d$ is log dividend growth and $\rho$ is a constant slightly below one. In words: if prices are high relative to dividends today, then either future dividends must grow fast, or future returns must be low, or the ratio must stay high forever (Figure~\ref{fig:07-identity}). The identity does not say which. Whether returns are predictable is the question of which side of the identity does the work.

\providecolor{fenavy}{HTML}{1B2A6B}

\begin{figure}[!htbp]
  \centering
  \begin{tikzpicture}[
    declare function={
      normal(\x)   = 4.72 - 3.3*exp(-\x/2.9);
      inverted(\x) = 3.8 + 1.35*exp(-\x/2.5);
    }]
  \begin{axis}[
    width=0.43\linewidth, height=0.27\linewidth,
    scale only axis,
    axis lines=left,
    axis line style={-stealth, thin},
    xmin=0, xmax=30.5, ymin=0, ymax=6.2,
    xtick={0,2,5,10,20,30},
    ytick={0,2,4,6},
    xlabel={Maturity (years)},
    ylabel={Yield (\%)},
    samples=150, domain=0.3:30,
    clip=false,
  ]
    \addplot[fenavy, very thick] {normal(x)};
    \addplot[thick, dashed, gray!30!black] {inverted(x)};

    % Term spread: 10-year yield minus short rate, on the normal curve
    \draw[fenavy, thick, <->] (axis cs:10.5,{normal(0.3)}) -- (axis cs:10.5,{normal(10)});
    \node[fenavy, anchor=north west, align=left] at (axis cs:11,3.6)
      {term spread:\\10-year minus short rate};

    % Curve labels
    \node[fenavy, anchor=west] at (axis cs:30.3,{normal(30)}) {normal (upward sloping)};
    \node[gray!30!black, anchor=west] at (axis cs:30.3,{inverted(30)}) {inverted};
  \end{axis}
  \end{tikzpicture}
  \caption{Two shapes of the yield curve (illustrative). Normally, long-term yields exceed short-term ones and the term spread is positive. Before recessions the curve often flattens or inverts.}
  \label{fig:07-curve}
\end{figure}

\providecolor{fenavy}{HTML}{1B2A6B}

\begin{figure}[!htbp]
  \centering
  \scalebox{0.8}{%
\begin{tikzpicture}[
    csbox/.style={draw=fenavy, thick, rounded corners=2pt, align=center,
                  minimum width=3.6cm, minimum height=1.15cm, inner sep=4pt},
    csarr/.style={-stealth, thick},
  ]
    % Top: today's price relative to dividends
    \node[csbox, fill=fenavy!12, minimum width=4cm, minimum height=1.35cm] (cstop) at (0,0)
      {Prices high relative to\\dividends today};

    % Three outcomes
    \node[csbox] (csret) at (-4.3,-2.9) {future returns\\are low};
    \node[csbox] (csdiv) at (0,-2.9)    {future dividends\\grow fast};
    \node[csbox] (csbub) at (4.3,-2.9)  {the ratio stays\\high forever};

    \draw[csarr] (cstop.south) ++(-1,0) -- (csret.north);
    \draw[csarr] (cstop.south) -- (csdiv.north);
    \draw[csarr] (cstop.south) ++(1,0) -- (csbub.north);

    \node[fenavy, font=\small] at (-5.1,-1.45) {returns are predictable};

    \node[font=\small] at (0,-4.3)
      {at least one must happen; the identity does not say which};
  \end{tikzpicture}%
}
  \caption{The Campbell--Shiller identity in words. A high price relative to
    dividends must be followed by low returns, fast dividend growth or a
    permanently high ratio. Return predictability is the first branch.}
  \label{fig:07-identity}
\end{figure}

\textbf{Two readings of predictability.} If expected returns vary over time, there are two broad explanations:

\begin{itemize}
  \item \textbf{Rational time-varying risk premia.} In bad times investors are poorer, more \textbf{risk-averse} (less willing to bear risk) or more exposed to risk, so they demand higher expected returns; prices fall until expected returns are high enough. Predictability is then compensation, not a free lunch (Fama and French 1989; Cochrane 2008).
  \item \textbf{Behavioral mispricing.} Investors \textbf{extrapolate} recent good news, assuming it will continue; prices overshoot, and later returns are low as they correct. Predictability is then a sign of mispricing (Shiller 2000).
\end{itemize}

Both readings imply that high valuations precede low returns. They differ on whether an investor can profit by acting on that, which is the question this chapter tests.

\textbf{Predictive regressions.} The standard test regresses the return over the next $k$ months on today's predictor, as in Chapter~\ref{ch:05}. Because valuation ratios move slowly, regressions at long horizons use overlapping returns: each monthly observation of a 10-year return shares 119 months with the next. Chapter~\ref{ch:03} showed how overlap inflates ordinary t-statistics; this chapter uses the corrections developed there.

\section{The claim}
\label{sec:07-claim}

Stated so that it can fail:

\begin{callout}
\textbf{Claim.} Valuation ratios and macroeconomic variables forecast the U.S. equity premium out of sample better than the historical mean benchmark.
\end{callout}

The last clause is what makes the claim testable. A pattern visible in the full history, with coefficients estimated using data the investor did not have, is not enough. The claim is about forecasts that could have been made in real time.

The chapter tests four predictors, chosen because each has a large literature and a clear economic story:

\begin{center}
\small
\begin{tabularx}{0.9\linewidth}{@{}>{\raggedright\arraybackslash}p{0.30\linewidth}>{\raggedright\arraybackslash}p{0.20\linewidth}>{\raggedright\arraybackslash}X@{}}
\toprule
Predictor & Expected sign & Story \\
\midrule
Dividend--price ratio & High ratio, high returns & Cheap stocks pay more per dollar invested \\
CAPE (entered with its sign reversed) & High CAPE, low returns & Prices high relative to normal earnings must eventually fall back, or earnings rise \\
Term spread (long-term minus short-term yield) & Wide spread, high returns & The spread is high near business-cycle troughs, when risk premia are high (Fama and French 1989) \\
The ``Fed model'': earnings yield minus the 10-year Treasury yield & High spread, high returns & Stocks are cheap when their earnings yield exceeds bond yields \\
\bottomrule
\end{tabularx}
\end{center}

The claim is tested at three horizons: one month, one year and ten years. They correspond to three uses: tactical trading, annual asset allocation, and long-term planning.

\section{The original evidence}
\label{sec:07-evidence}

The case for predictability was built in three waves, and the case against it in a fourth.

\textbf{Valuation ratios (1987--1998).} Fama and French (1988) regressed U.S. stock returns on dividend yields and found that the explanatory power rose with the horizon: dividend yields explained little of monthly returns but often more than a quarter of the variance of returns over two to four years. Campbell and Shiller (1988) derived the present-value identity of Section~\ref{sec:07-how} and showed that dividend--price ratios and smoothed earnings forecast returns. Campbell and Shiller (1998), writing as U.S. valuations reached record highs, used ten-year average earnings, the measure now called CAPE, to argue that the market's expected returns over the following decade were unusually low. Shiller had presented similar evidence to the Federal Reserve Board in December 1996, days before Alan Greenspan's ``irrational exuberance'' speech.

\textbf{Business-cycle variables (1986--1989).} Keim and Stambaugh (1986), Campbell (1987) and Fama and French (1989) found that interest-rate spreads also predicted returns: notably the term spread between long- and short-term yields, and the \textbf{default spread} between the yields of low- and high-grade corporate bonds. Fama and French linked the pattern to the business cycle: expected returns were high near \textbf{troughs}, the low points of recessions, when spreads were wide and investors' willingness to bear risk was low. Their interpretation made predictability compatible with efficient markets, as time-varying compensation for risk.

\textbf{A consensus forms.} By the early 2000s, predictability from valuation ratios was widely taught as a stylized fact, and Cochrane's (2011) presidential address to the American Finance Association built a research agenda on the premise that discount rates, not expected cash flows, account for most of the movement in valuations.

\textbf{The challenge (1999--present).} Statisticians had warned that the evidence was fragile.

\begin{itemize}
  \item Stambaugh (1999) showed that persistent predictors whose surprises are correlated with returns bias slopes and t-statistics in finite samples.
  \item Valkanov (2003) and Boudoukh, Richardson and Whitelaw (2008) showed that long-horizon regressions add little information beyond short-horizon ones: the rising $R^2$ is largely a mechanical consequence of overlapping data and persistent predictors.
  \item Ang and Bekaert (2007), using data from several countries, found that long-horizon predictability from dividend yields largely disappeared once standard errors were corrected, while dividend yields helped at short horizons in combination with short-term interest rates.
  \item Goyal and Welch (2008) showed that almost none of the published predictors beat the historical mean in walk-forward forecasts, and Goyal, Welch and Zafirov (2024) found the same for predictors proposed since.
\end{itemize}

The defense did not concede. Campbell and Thompson (2008) showed that simple economic restrictions improve out-of-sample performance and that tiny $R^2$ values can be economically valuable. Cochrane (2008) argued that the absence of dividend-growth predictability implies return predictability through the present-value identity, so the evidence should be judged jointly rather than one regression at a time.

\section{Lab~\thechapter: replicating Campbell and Shiller}
\label{sec:07-lab}

\begin{callout}
\textbf{Lab~\ref{ch:07} at a glance.} \emph{Claim on trial\index{claims-on-trial method}:} valuation ratios and macroeconomic variables forecast stock returns better than the historical mean benchmark. \emph{Data:} the Goyal--Welch monthly data from 1927, the same file used in Lab~\ref{ch:05}, which replicates Goyal and Welch (2008), which includes the S\&P 500 index, twelve-month dividends and earnings, Treasury yields, inflation and CRSP total returns. CAPE is built from the index, earnings and inflation; it begins in 1923, once ten years of inflation-adjusted earnings are available. Optionally, Robert Shiller's long-run file extends CAPE and real returns back to 1881 for the ten-year analysis. \emph{Code:} repository folder \texttt{ch07-predictability/}, with a guided notebook.
\end{callout}

\textbf{The setup.} Lab~\ref{ch:07} replicates the empirical evaluation of Campbell and Shiller (1988, 1998) and Goyal and Welch (2008) across four predictive models and three horizons:

\begin{itemize}
  \item \textbf{Part A, replication.} Regressions of future log excess returns on each predictor at each horizon, with three t-statistics: ordinary least squares, which ignores the overlap; Newey--West with lags equal to the horizon minus one; and Hodrick's (1992) standard errors, which avoid overlapping errors altogether. Lab Figure~\ref{ch:07}.1 shows the classic pattern to look for: $R^2$ rising with the horizon.
  \item \textbf{Part B, walk-forward.} Forecasts from January 1965 using only data available at each date, with training observations purged when their targets had not yet been observed (Chapter~\ref{ch:05}). Out-of-sample $R^2$ against the historical mean, the Clark--West test, and the Campbell--Thompson restrictions. Lab Figure~\ref{ch:07}.2 sets these beside the in-sample results.
  \item \textbf{Part C, stress tests.} Section~\ref{sec:07-stress} describes these.
  \item \textbf{Part D, ranges.} Annualized ten-year real returns regressed on starting CAPE, with the walk-forward forecast error used to set an honest range around the forecast for today's CAPE (Lab Figure~\ref{ch:07}.5).
\end{itemize}

\textbf{What readers will find.} The lab reports every number from the data the reader supplies, so the results update as the data do. The pattern to test is the one the literature reports: in-sample $R^2$ of around 1\% at one month, larger at one year, and large at ten years; out-of-sample $R^2$ near or below zero at the short horizons for most predictors; and a ten-year relationship with CAPE that holds out of sample but leaves a range several percentage points wide. When the lab was run on simulated data in which the dividend--price ratio has weak true predictive power, the bootstrap of Part C showed why the published t-statistics overstate the evidence: the 95\% critical value under the null was about 2.6 rather than the textbook 1.645.

\section{Stress tests}
\label{sec:07-stress}

Four stress tests apply the earlier chapters' tools to the claim.

\textbf{1. Overlapping data (Chapter~\ref{ch:03}).} At the ten-year horizon, consecutive monthly observations share 119 of their 120 months. A century of data contains only about ten non-overlapping decades. Ordinary t-statistics treat the 1,000-plus overlapping observations as if they were independent and can overstate significance many times over. The lab reports OLS, Newey--West and Hodrick t-statistics side by side; the gap between the first and the others is the size of the overlap problem.

\textbf{2. The Stambaugh bias.} Valuation ratios are extremely persistent, and their surprises are strongly negatively correlated with return surprises: when prices jump, the dividend--price ratio falls. In finite samples this biases the estimated slope upward and makes t-statistics too large (Stambaugh 1999; Nelson and Kim 1993). The lab measures the bias directly with a bootstrap: it keeps each predictor's persistence and the correlation of its surprises with returns, but rebuilds returns with no predictability at all, and runs the same regression on 500 artificial histories. Lab Figure~\ref{ch:07}.4 shows the resulting distribution of t-statistics under the null, which is shifted to the right of zero, and compares the actual t-statistic with that distribution rather than with the textbook normal.

\textbf{3. Different decades.} A relationship that reflects a stable economic mechanism should appear in different periods. The lab estimates each predictor's one-year slope in three subperiods (1927--1964, 1965--1994 and 1995 onward) and over rolling 20-year windows (Lab Figure~\ref{ch:07}.3). The dividend--price ratio is the classic test case: the ratio fell to record lows in the late 1990s while returns stayed high, and many studies find that its slope and significance depend heavily on whether that period is included.

\textbf{4. Out of sample and with restrictions (Chapter~\ref{ch:05}).} Walk-forward forecasts with purging, the Clark--West test, and the Campbell--Thompson restrictions. These are the decisive tests for the claim as stated, because the claim is about forecasts an investor could have made.

\textbf{What the stress tests do not cover.} All four use U.S. data. International evidence is weaker and more mixed, which matters because the U.S. market's history is one of the most successful on record, a selection problem Chapter~\ref{ch:02} discussed. The open questions in Section~\ref{sec:07-open} return to this.

\section{Verdict}
\label{sec:07-verdict}

The verdict depends on the horizon and on what the predictions are for. Based on the published evidence through the Goyal--Welch--Zafirov sample, and to be confirmed by readers' runs of Lab~\ref{ch:07} on data through 2025:

\begin{center}
\small
\begin{tabularx}{0.9\linewidth}{@{}>{\raggedright\arraybackslash}p{0.21\linewidth}>{\raggedright\arraybackslash}p{0.11\linewidth}>{\raggedright\arraybackslash}p{0.13\linewidth}>{\raggedright\arraybackslash}X@{}}
\toprule
Use & Horizon & Verdict & Why \\
\midrule
Market timing: moving in and out of stocks & 1 month to 1 year & \textbf{Fails} & In-sample t-statistics shrink after correcting for overlap and the Stambaugh bias; walk-forward forecasts rarely beat the historical mean, and the predictors that help do so only in some decades \\
Long-run planning: setting expected-return assumptions & 10 years & \textbf{Survives, weakly} & High CAPE has been followed by lower ten-year returns, a pattern that has generally held up in walk-forward tests, but with roughly ten independent decades of evidence the relationship is imprecise and the honest range is several percentage points wide \\
The Fed model as a valuation signal & Any & \textbf{Fails} & Comparing a real earnings yield with a nominal bond yield has no sound economic basis, and it has not forecast returns \\
\bottomrule
\end{tabularx}
\end{center}

The verdict on the claim as stated, that these variables forecast the equity premium well enough to beat a constant expected return in real time, is therefore \textbf{fails for market timing}, with a narrower claim, that valuations shift the distribution of decade-long returns, surviving weakly. The verdict is dated by the sample tested and should be revisited as data accumulate: a decade of new returns adds one new independent observation to the ten-year evidence.

\section{What investors can use}
\label{sec:07-investors}

The useful lesson from a century of debate is not a timing signal but a planning input.

\begin{itemize}
  \item \textbf{Set long-run return assumptions with valuations in mind.} A retirement plan, an endowment's spending rule or a pension's discount rate that assumes the historical average return regardless of starting valuations ignores the most robust part of this evidence. When CAPE is high, plan for lower decade-long returns; when it is low, higher ones.
  \item \textbf{Plan for the range, not the point.} Lab~\ref{ch:07}'s Part D reports a 90\% range for the next decade's annualized real return. The range is wide, typically several percentage points either side of the central forecast, and the plan should work across it.
  \item \textbf{Do not trade on it.} Valuations have been high, or low, for years at a time before returns responded. An investor who reduced equities whenever CAPE looked expensive would often have spent long periods underinvested while the market rose.
  \item \textbf{Distrust precise forecasts.} Any forecast of next year's market return stated to a decimal place, or a valuation ``fair value'' target for the index, implies a precision this evidence does not support.
\end{itemize}

\textbf{Questions to ask a fund manager or adviser who times the market with valuations:}

\begin{enumerate}
  \item Were your signals and their thresholds chosen before the period you are showing me? What did you try and discard?
  \item How did the strategy perform in real time, after it was designed, not in a backtest?
  \item What t-statistics do you get after correcting for overlapping returns and the persistence of your predictor?
  \item How long has the signal been wrong in the past, and could I bear being out of the market for that long?
  \item What is the strategy's out-of-sample $R^2$ against simply assuming the average return?
\end{enumerate}

\begin{callout}
\textbf{What Failed: the Fed model.} In July 1997 the Federal Reserve's Monetary Policy Report to Congress included a chart comparing the stock market's earnings yield with the ten-year Treasury yield. The strategist Ed Yardeni named the comparison the ``Fed model'': stocks are cheap when their earnings yield is above bond yields and expensive when below. The Fed never endorsed it, but it became one of Wall Street's most quoted valuation tools. Asness (2003) argued that it rests on a confusion: earnings yields are real quantities, because earnings grow with inflation, while bond yields are nominal, so comparing them builds in \textbf{money illusion}, the mistake of treating nominal and real amounts as the same, of the kind Modigliani and Cohn (1979) described. Its empirical record matches the logic. Asness found that it described how investors set prices but not what returns followed, and it had little power to forecast long-run returns. The relationship it relies on was also unstable: from the late 1990s, bond yields and earnings yields moved in opposite directions for long periods. The claim that the Fed model identifies when stocks are cheap \textbf{fails}. Lab~\ref{ch:07} includes it as a fourth predictor so readers can check.
\end{callout}

\section{Open questions}
\label{sec:07-open}

\begin{itemize}
  \item \textbf{International evidence.} Do valuation ratios predict long-run returns outside the United States, using long histories for many countries? The answer bears on whether U.S. evidence reflects a mechanism or one fortunate market.
  \item \textbf{Changes in payout policy.} Share buybacks have replaced much of the dividend payout since the 1980s, which may have shifted the dividend--price ratio's level permanently. How should valuation ratios be adjusted, and does adjustment restore their predictive power?
  \item \textbf{Accounting changes in earnings.} Changes in accounting rules, such as write-down rules for goodwill, affect reported earnings and therefore CAPE. How large are these effects, and how should a long-run series correct for them?
  \item \textbf{Real-time data.} Lab~\ref{ch:07} uses earnings as dated in the data file, not as first reported. How much does using genuinely real-time earnings change the evidence?
  \item \textbf{Combining signals.} Do forecast combinations, Bayesian methods that impose economic priors, or machine-learning models (Chapter~\ref{ch:21}) deliver forecasts that survive the stress tests of this chapter?
\end{itemize}

\section*{Exercises}
\addcontentsline{toc}{section}{Exercises}

\subsection*{Conceptual (all readers)}

\begin{enumerate}
  \item Use the present-value identity of Section~\ref{sec:07-how} to explain why a high dividend--price ratio must be followed by high returns, high dividend growth, or a permanently high ratio. Which of these does the evidence favor?
  \item Explain why the $R^2$ of a regression of 10-year returns on a persistent predictor can be large even if the predictor has little information about 1-month returns.
  \item Why does the Stambaugh bias push t-statistics up for the dividend--price ratio, and why would it push them the other way for a predictor whose surprises were positively correlated with returns?
  \item An adviser says CAPE predicts that the next ten years will deliver 3\% a year after inflation. What information is missing from that statement?
\end{enumerate}

\subsection*{Lab extensions (technical readers)}

\begin{enumerate}
  \setcounter{enumi}{4}
  \item Re-run Lab~\ref{ch:07} ending in 1990 (\texttt{-{}-end 1990-12}), roughly the data available to the original authors. Then run it from 1991. How different are the conclusions?
  \item Compare the OLS, Newey--West and Hodrick t-statistics for the 10-year horizon. By what factor does ignoring overlap overstate significance?
  \item Lag CAPE by four months to allow for delays in reported earnings. Does the evidence change?
  \item Supply Shiller's data and repeat Part D from 1881. How does the longer history change the slope and the width of the range?
  \item Obtain long-run data for another country's stock market and repeat Part D. Does the CAPE relationship hold?
  \item Build a combination forecast of the four predictors, as in Rapach, Strauss and Zhou (2010), and evaluate it with the stress tests of Section~\ref{sec:07-stress}.
\end{enumerate}